\documentclass[12pt]{article}
\usepackage{bbm}
\usepackage{fancyhdr}
\usepackage{titling}
\usepackage{enumitem}
\usepackage{amsmath}
\usepackage{dsfont}
\usepackage{amsfonts}
\usepackage{amssymb}
\usepackage{amsthm}
\usepackage{geometry}
\usepackage{tikz}
\usepackage{algorithm}
\usepackage{algpseudocode}
\geometry{verbose,tmargin=1in,bmargin=1in,lmargin=1in,rmargin=1in}
\renewcommand{\vec}[1]{\mathbf{#1}}
\newcommand{\dec}{\text{dec}}
\newcommand{\divi}{\mathbin{\textnormal{div}}}
\title{Homework 1}
\author{Satvik Prasad}
\date{}
\newtheorem{lemma}{Lemma}[section]
\fancypagestyle{plain}{%
  \fancyhf{}
  \fancyhead[L]{\theauthor}
  \fancyhead[R]{\thedate}
}
\pagestyle{plain}
\makeatletter
\def\@maketitle{%
  \newpage
  \null
  \vskip 1em%
  \begin{center}%
    \let \footnote \thanks
    {\LARGE \@title \par}%
    \vskip 1em%
  \end{center}%
  \par
\vskip 1em}
\makeatother
\usepackage{cmbright}

\newcommand{\mat}[1]{\mathbf{#1}}

% ---------- Probability theory ----------
% Core operators (optional argument-free; use as \E{X}, \Var{X}, \Cov{X}{Y})
\DeclareMathOperator{\EOp}{\mathbb{E}}
\DeclareMathOperator{\VarOp}{Var}
\DeclareMathOperator{\CovOp}{Cov}
\DeclareMathOperator{\CorrOp}{Corr}

\renewcommand{\P}[1]{\mathbb{P}\left(#1\right)}        % \P{A}
\newcommand{\E}[1]{\EOp\left[#1\right]}                % \E{X}
\newcommand{\Var}[1]{\VarOp\left(#1\right)}            % \Var{X}
\newcommand{\Cov}[2]{\CovOp\left(#1,\,#2\right)}       % \Cov{X}{Y}
\newcommand{\Corr}[2]{\CorrOp\left(#1,\,#2\right)}     % \Corr{X}{Y}
\newcommand{\SD}[1]{\operatorname{SD}\left(#1\right)}  % \SD{X}

% Subscripted versions, e.g. \Esub{\theta}{X}, \Psub{\theta}{A}
\newcommand{\Esub}[2]{\EOp_{#1}\left[#2\right]}
\newcommand{\Psub}[2]{\mathbb{P}_{#1}\left(#2\right)}

% Conditioning, indicators, independence
\newcommand{\given}{\,\middle|\,}                      % \E{X \given Y}
\newcommand{\ind}[1]{\mathbbm{1}\left\{#1\right\}}     % \ind{X > 0}
\newcommand{\indep}{\perp\!\!\!\perp}                  % X \indep Y
\newcommand{\iid}{\overset{\text{iid}}{\sim}}          % X_i \iid \Normal(0,1)

% Convergence
\newcommand{\convp}{\xrightarrow{\;p\;}}               % in probability
\newcommand{\convd}{\xrightarrow{\;d\;}}               % in distribution
\newcommand{\convas}{\xrightarrow{\text{a.s.}}}        % almost surely
\newcommand{\convLp}[1]{\xrightarrow{L^{#1}}}          % in L^p

% Probability spaces
\newcommand{\F}{\mathcal{F}}                           % sigma-algebra
\newcommand{\B}{\mathcal{B}}                           % Borel sigma-algebra
\newcommand{\pspace}{(\Omega, \F, \mathbb{P})}

% Common distributions
\newcommand{\Normal}{\mathcal{N}}
\newcommand{\Bern}{\operatorname{Bernoulli}}
\newcommand{\Bin}{\operatorname{Binomial}}
\newcommand{\Pois}{\operatorname{Poisson}}
\newcommand{\Unif}{\operatorname{Uniform}}
\newcommand{\Exp}{\operatorname{Exp}}
\newcommand{\Geom}{\operatorname{Geometric}}
\newcommand{\Gam}{\operatorname{Gamma}}
\newcommand{\Beta}{\operatorname{Beta}}

% Number sets
\newcommand{\R}{\mathbb{R}}
\newcommand{\N}{\mathbb{N}}
\newcommand{\Z}{\mathbb{Z}}

\begin{document}

\begin{enumerate}
    \item \begin{enumerate}
        \item D reduces to all of \{C, F, E, G\}, meaning if D is NP-hard \{C, D, E, F, G\} must be NP-hard.

        \item If $P \neq NP,$ then no NP-hard problem can be in P, otherwise all of NP is in P. Thus, \{A, B, H\} can be in P.
    \end{enumerate}

    \item \begin{enumerate}
        \item To prove B is NP-complete, we need to prove it is NP-hard since B is already in NP. Hence, we need to show that all of NP reduces to B. We can do this by showing A reduces to B, since A is already NP-hard.

        \item If the reductions don't have a polynomial time complexity, then the statement $A \to B$ in $\omega(\text{poly}(n))$ doesn't necessarily imply $A$ is in P if B is in P, since the reduction itself takes longer than polynomial time.

        \item For there to be a reduction $A \to B,$ we want to be able to identify all solutions of $A$ by solving $B.$ Hence, if there is a YES-instance of inputA, running our algorithm on $B$ should produce a YES-instance of inputB, meaning all YES-instances on A are YES-instances on B. Conversely, we don't want to mis-identify any non-solutions, meaning if we get a YES-instance on inputB, it better be a YES-instance on inputA. Hence, we need a YES-instance on inputA iff we have a YES-instance on inputB.
    \end{enumerate}

    \item The decision "is there a path of length $\geq k$" is monotonic (non-decreasing) on $k.$ That is, if there is a path of length $\geq k,$ there is a path of length $\geq r$ for $r < k.$ Similarly, if there is no path of length $\geq k,$ there is no path for length $r > k,$ meaning there is some decision boundary.

    \begin{itemize}
        \item Assume we have $\texttt{DecideLongestPath}(G, k)$ that runs in $T(|V|, |E|)$ time. Then, we can binary search on $k,$ with the update rule $r \gets \texttt{mid} - 1$ if NO, and $l \gets \texttt{mid} + 1$ if YES. The final answer will be $k_\text{longest} = l - 1$. 
        
        Our initial range on $k$ is $[0, |V| - 1],$ since any path with $> n -1$ edges must be incident on $> n$ vertices, and we cannot re-visit vertices. Since this requires $O( \log |V|)$ iterations of $\texttt{DecideLongestPath},$ this algorithm runs in 
        $$
            \mathcal O(T(|V|, |E|) \cdot \log |V|)
        $$
        time.

        \item First, we run the above algorithm to determine the length of the longest path, call this $L.$ Then, we can iterate through every edge of the graph, remove it from $G,$ and run $\texttt{DecideLongestPath}.$ If the result $= L,$ we know there is a longest path that does not take the edge, and we can remove this. We continue this process until we have $L$ edges left, at which point we can BFS from a vertex along this path to reach an endpoint, and BFS from there.

        \paragraph{Justification} We prove that every edge only needs to be iterated a single time. Consider a step of the iteration over edge $e.$ We have two cases:
        \begin{itemize}
            \item Edge $e$ is not taken on any longest path in this graph, or there are multiple such longest paths. In this case, we permanently remove $e$ from $G,$ and it won't recur in our iteration.

            \item There is a single longest path in what remains of $G$ (after modifications), and edge $e$ is along this path. In this case, we know we must keep $e$ since it is along our only remaining longest path (that we wish to detect), and we don't need to iterate over it again.
        \end{itemize}

        Hence, we can find the longest path in 
        $$
        \mathcal O(T(|V|, |E|) \cdot \log |V| + T(|V|, |E|) \cdot |E| + |V| + |E|) = \mathcal O(T(|V|, |E|) \cdot (|E| + \log |E|))
        $$
        time since $\mathcal O(\log |V|) = \mathcal O(\log |E|),$ and $\mathcal O(|V|) = \mathcal O(|E|)$ as required.
    \end{itemize}

    \item \paragraph{The bounded spanning tree problem (BST) is in NP}
    \begin{algorithm}
        \caption{Verify spanning-tree}
        \begin{algorithmic}
            \Procedure{VerifySpanningTree}{$G, k, T$}
                \For{$v \in G$}
                    \If{$v \not \in T$} \Comment{Check if spanning, $\mathcal O(|V|^2)$ naive}
                        \State \Return \texttt{false}
                    \EndIf
                \EndFor

                \vspace{3mm}

                \State $\texttt{visited},\, \texttt{has\_cycle} \gets \texttt{dfs}(T)$ \Comment{Check if tree, $\mathcal O(|V| + |E|)$}
                \If{$\texttt{has\_cycle} \lor (\texttt{visited} \neq V(T))$}
                    \State \Return \texttt{false}
                \EndIf

                \vspace{3mm}
            
                \For{$v \in T$} \Comment{Check if $\deg$ valid, $\mathcal O(|V|)$}
                    \If{$\deg v > k$}
                        \State \Return \texttt{false}
                    \EndIf
                \EndFor

                \State \Return \texttt{true}
            \EndProcedure
        \end{algorithmic}
    \end{algorithm}
    Such an algorithm verifies that the tree is an spanning tree, and has bounded degree in polynomial time, meaning BST $\in$ NP.

    \paragraph{BST is NP hard} We reduce HamPath to BST. Notice that any tree $T$ where every vertex $v$ has degree $\leq 2$ must be a path:
    \begin{itemize}
        \item There is a vertex $v \in T$ with degree $0.$ Since this is a tree, there are no other vertices in the tree (otherwise this vertex would be disconnected and $T$ no longer a tree), and the tree is a lone vertex. In this case, $T$ is a path.

        \item $T$ is not a lone vertex. Then, $T$ has vertices with degree $1$ or $2.$ We continue via induction.

        Assume all trees $T$ with $n = k \geq 2$ vertices and all vertices with degree $\leq 2$  are paths. Then, consider an arbitrary tree $T$ with $k + 1$ vertices and degree $2.$ 
        
        There must be a vertex in $T$ with degree $1.$ Otherwise, we can construct a cycle by starting at some start vertex, and taking the vertex that we are not incident from. Since there are $k + 1$ vertices this process must eventually repeat a vertex, at which point we have a cycle.

        Hence, take this vertex $u$ with degree $1$ and remove it. We are left with a graph with $\geq 2$ vertices with degree $2,$ meaning it is a path. The vertex incident from $u$, $u^\star,$ previously had degree $2$ (otherwise if it had degree $1$, would be $P_2$), and hence by our inductive hypothesis $T \setminus \{u\}$ is a path with endpoint $u^\star.$ Then, $T$ adds edge $\{u, u^\star\},$ and since $u^\star$ was an endpoint, $T$ is still a path, concluding the inductive step.

        For our base case, the only tree with $2$ vertices is $P_2,$ which is a path
    \end{itemize}

    In all cases, $T$ with degree bounded by 2 is a path. Hence, if $T$ is spanning and degree bounded by 2, $T$ is a path incident on all vertices, meaning $T$ is a Hamiltonian path. Conversely, if $P$ is a Hamiltonian path, every vertex must have degree $\leq 2,$ and the path must be spanning, meaning $P$ is also a spanning tree with degree $\leq 2.$

    Thus, if Hamiltonian path exists in $G,$ then a spanning tree with degree $k\leq 2$ also exists, and vice versa.

    Therefore, we can solve HamPath by simply running BST on our graph $G$ with $k = 2.$ The reduction is $\mathcal O(|V|, |E|).$ By definition, since HamPath is NP-hard, and it reduces to BST, all problems in NP must reduce to BST, meaning BST is NP-hard.

    Hence, BST is NP-hard and in NP, i.e., BST is NP complete.
\end{enumerate}
\end{document}